\chapter{Diretrizes Metodológicas do Ensino}
\label{chp:diretrizes_pedagogicas}

O \ppccurso da Universidade Federal de Uberlândia adota uma concepção pedagógica centrada na formação integral do(a) estudante, combinando o desenvolvimento de conhecimentos técnico-científicos com habilidades práticas, atitudes éticas e capacidades reflexivas. As diretrizes metodológicas são pautadas por uma abordagem orientada por competências, conforme as Diretrizes Curriculares Nacionais do Curso de Graduação em Engenharia~\cite{cne_ces_2_2019,cne_ces_1_2021}, valorizando o protagonismo discente, a interdisciplinaridade e a articulação entre teoria e prática profissional. A referência anterior à Resolução CNE/CP nº 1/2021, que trata da Educação Profissional e Tecnológica, foi removida por não se aplicar a este Bacharelado em Engenharia.

\section{Princípios Metodológicos}

As práticas de ensino-aprendizagem são orientadas pelos seguintes princípios:

\begin{itemize}
    \item Integração entre teoria e prática, por meio de aulas experimentais, projetos, estudos de caso, práticas laboratoriais e atividades extensionistas;
    \item Interdisciplinaridade, promovendo conexões entre eletricidade, eletrônica, controle, automação, redes industriais e ciências humanas e sociais aplicadas;
    \item Flexibilidade curricular, com espaço para disciplinas optativas e reconhecimento de saberes adquiridos em experiências acadêmicas e profissionais;
    \item Aprendizagem ativa, com incentivo ao uso de metodologias como sala de aula invertida\footnote{\textit{Flipped Classroom}, modelo pedagógico em que a lição de casa tradicional se torna o aprendizado em casa, e as atividades em sala de aula são focadas em exercícios e discussões mais aprofundadas.}, aprendizagem baseada em projetos (PBL), resolução de problemas (ABP), estudos dirigidos e simulações computacionais de processos industriais;
    \item Tecnologias digitais na educação, aproveitando plataformas virtuais, ambientes de simulação, ferramentas colaborativas e sistemas de avaliação digital;
    \item Educação híbrida, integrando atividades presenciais e a distância de forma planejada, conforme a natureza dos componentes curriculares e as necessidades dos(as) estudantes;
    \item Formação continuada, com incentivo à aprendizagem ao longo da vida, à autonomia técnica e à atualização diante da rápida evolução das tecnologias de automação.
\end{itemize}

\section{Integração Pedagógica na Área Básica de Ingresso (ABI)}
\label{sec:integracao_pedagogica_abi}

O \ppccurso é organizado na ABI com o Curso Superior de Tecnologia em Automação Industrial (\autoref{chp:abi}). Do ponto de vista pedagógico, essa organização estrutura-se em duas fases articuladas entre si:

\begin{itemize}
    \item nos cinco primeiros períodos, comuns às duas graduações, os componentes curriculares são planejados para desenvolver competências e conteúdos de fundamentação — científica, tecnológica e humanística — igualmente relevantes para a trajetória tecnológica e para a trajetória em Engenharia, sem antecipar conteúdos exclusivos de nenhuma das duas;
    \item a partir do sexto período, a formação progride para a trajetória específica escolhida pelo estudante, aprofundando os conteúdos e as competências próprias de cada perfil profissional (\autoref{chp:perfil_egresso}) — sem retomar, como pré-requisito pedagógico adicional, conteúdos já trabalhados na etapa comum.
\end{itemize}

Essa progressão preserva os perfis profissionais próprios de cada graduação: a etapa comum não dilui as competências específicas do \ppccurso nem as do Curso Superior de Tecnologia em Automação Industrial, apenas antecipa, de forma pedagogicamente coerente, os conteúdos que ambas as formações efetivamente compartilham. Antes da escolha ao final do quinto período, o estudante recebe orientação acadêmica sobre as duas trajetórias, seus requisitos e suas implicações formativas, de modo a subsidiar uma decisão informada (\autoref{sec:orientacao_academica_abi}).

\section{Uso de Educação a Distância (EaD)}
\label{sec:ead}

O \ppccurso da UFU adota a Educação a Distância (EaD) como uma ferramenta complementar de ensino, com o objetivo de promover metodologias ativas, acesso ampliado a conteúdos, flexibilidade pedagógica e inovação tecnológica. Essa prática está em conformidade com o Decreto nº 12.456, de 19 de maio de 2025~\cite{mec_decreto_ead_12456_2025}, e com a Portaria MEC nº 378, de 19 de maio de 2025~\cite{mec_portaria_378_2025}.

Em cursos de graduação presenciais, como é o caso deste, é permitida a inclusão de até \textbf{30\% da carga horária total} do curso com atividades síncronas ou assíncronas realizadas por EaD, conforme o Art. 10 do referido Decreto — percentual que corresponde exatamente ao configurado para este curso (\texttt{curriculo.percentual\_maximo\_ead}, aba \texttt{Perfil} da matriz curricular).

\subsection{Definições e diretrizes fundamentais}

As modalidades de EaD previstas neste PPC incluem:

\begin{itemize}
    \item Atividades síncronas: realizadas com a presença simultânea de docentes e discentes por meio de recursos de áudio e vídeo (ex: videochamadas, salas virtuais ao vivo);
    \item Atividades assíncronas: realizadas em tempos distintos por docentes e discentes (ex: fóruns, videoaulas, tarefas em AVA, sala de aula invertida);
    \item Atividades síncronas mediadas: síncronas com grupo de até 70 estudantes por docente ou mediador, com controle de frequência obrigatório;
    \item Atividades presenciais obrigatórias: que podem ocorrer na sede da UFU, em laboratórios, em ambientes profissionais ou espaços de extensão, conforme o plano pedagógico.
\end{itemize}

\subsection{Planejamento e supervisão}
\label{sec:planejamento_ead}

Toda unidade curricular com carga EaD será planejada e supervisionada por docente responsável (professor regente), podendo haver apoio de mediadores pedagógicos e tutores, conforme previsão legal. Os materiais didáticos seguirão os princípios de:
\begin{itemize}
    \item alinhamento com os objetivos de aprendizagem e as competências da disciplina;
    \item acessibilidade e diversidade de recursos;
    \item compatibilidade com as Diretrizes Curriculares Nacionais e os referenciais da UFU.
\end{itemize}

As atividades EaD do curso serão desenvolvidas preferencialmente nas plataformas institucionais Moodle UFU ou Microsoft Teams, que oferecem suporte ao acompanhamento das atividades síncronas e assíncronas, registro acadêmico, interação pedagógica e gestão de conteúdos digitais.

Componentes presenciais poderão incluir carga horária mediada por tecnologia quando essa previsão constar da matriz, da ficha do componente e do plano de ensino, observados os requisitos acadêmicos e institucionais. A soma dessas atividades deve permanecer dentro do limite de 30\% da carga horária total do curso, conforme o Decreto nº 12.456/2025; este PPC não estabelece um limite autônomo de 20\% por disciplina. As atividades poderão incluir tarefas digitais, videoaulas, fóruns, leituras e encontros síncronos, conforme a caracterização registrada para o componente.

\subsection{Avaliação das atividades EaD}

As avaliações de unidades curriculares com EaD seguirão os princípios gerais do curso, devendo observar:

\begin{itemize}
    \item obrigatoriedade de avaliações presenciais para atividades EaD que representem fração significativa da carga horária;
    \item desenvolvimento de habilidades de análise e síntese discursiva, compondo ao menos 1/3 da nota, salvo exceções práticas (Art. 23, §1º, III)~\cite{mec_decreto_ead_12456_2025};
    \item garantia da autenticidade e identificação dos(as) estudantes nas avaliações síncronas e presenciais.
\end{itemize}

\subsection{Compromisso institucional}

A UFU, como Instituição de Educação Superior pública federal, está automaticamente credenciada para a oferta complementar de atividades EaD, conforme Art. 15 do Decreto nº 12.456/2025. O \ppccurso reafirma seu compromisso com:

\begin{itemize}
    \item o uso responsável, pedagógico e inovador das tecnologias digitais;
    \item o respeito à legislação vigente, aos limites estabelecidos e à identidade do curso;
    \item a valorização da interação humana, do engajamento ativo e do aprendizado de excelência.
\end{itemize}

\textbf{Pendência institucional:} a matriz registra \path{oferta.status_validacao_institucional=pendente}. A oferta de carga mediada por tecnologia somente deve ser implantada após verificação das normas e dos procedimentos vigentes da UFU, além da observância do Decreto nº 12.456/2025 e da Portaria MEC nº 378/2025.

\section{Integração com Pesquisa Aplicada, Extensão e Vivência Profissional}

A metodologia do curso articula ensino, pesquisa aplicada e extensão como dimensões indissociáveis da formação universitária. Os(as) estudantes são incentivados(as) a participar de projetos de iniciação científica e tecnológica, das Atividades Curriculares de Extensão (Seção~\ref{sec:ace}), de atividades em laboratórios especializados, de estágios e de programas de qualificação complementar.

Além disso, o curso valoriza as aprendizagens desenvolvidas em contextos extracurriculares -- incluindo a mobilidade acadêmica nacional e internacional e a realização de componentes curriculares em outras instituições de ensino superior (Seção~\ref{sec:mobilidade}) --, reconhecendo-as, nos termos das normas institucionais vigentes da UFU, como parte legítima da formação do(a) estudante.

\section{Avaliação do Processo de Ensino-Aprendizagem}

A avaliação no curso é formativa, diagnóstica e contínua, buscando não apenas aferir resultados, mas orientar o desenvolvimento progressivo das competências previstas. As estratégias de avaliação são variadas e contextualizadas, incluindo:

\begin{enumerate}
    \item provas escritas e orais;
    \item projetos e trabalhos em grupo;
    \item atividades práticas e laboratoriais;
    \item autoavaliações e coavaliações;
    \item apresentações técnicas e relatórios;
    \item portfólios e produções técnicas autorais.
\end{enumerate}

A coerência entre métodos de ensino, objetivos da disciplina, competências visadas e formas de avaliação é elemento fundamental da qualidade do processo formativo.

\section{Uso Ético e Consciente de IA no Ensino-Aprendizagem}
\label{sec:IA-UFU}

O \ppccurso da UFU reconhece a Inteligência Artificial (IA), em especial a IA Generativa, como uma tecnologia transformadora com profundo impacto no ensino, na pesquisa aplicada e na prática profissional da automação industrial — área na qual este curso adota, inclusive, ênfase curricular específica (\ppcenfasecurricular). Alinhado ao Guia de Princípios para o Uso Ético e Responsável de IA da UFU~\cite{ufu_guia_ia_2025} e às discussões contemporâneas sobre o tema, este PPC não busca restringir o uso dessas ferramentas, mas sim promover uma cultura de letramento em IA, capacitando os(as) discentes a se tornarem curadores conscientes da cognição, em vez de meros consumidores de tecnologia.

Nossa abordagem se fundamenta na premissa de que o verdadeiro risco da IA não é a substituição do intelecto humano, mas a abdicação voluntária de nossa capacidade de pensar criticamente. Portanto, as seguintes diretrizes orientam o uso de IA no percurso formativo:

\subsection{Princípios Éticos e Integridade Acadêmica}

\begin{enumerate}
    \item \textbf{Supervisão Humana e Responsabilidade Final:} Toda e qualquer utilização de IA deve ser rigorosamente supervisionada. O(a) discente é o(a) autor(a) e responsável final por qualquer conteúdo que apresente, cabendo-lhe a verificação da acurácia, a validação das fontes e a originalidade do trabalho. A IA é uma ferramenta de apoio, não um substituto para o discernimento humano.

    \item \textbf{A Distinção entre Plausibilidade e Verdade:} Os modelos de IA Generativa são projetados para sintetizar conteúdo plausível, que se assemelha a padrões aprendidos, mas não necessariamente verdadeiro ou factualmente correto. É imperativo que os(as) estudantes compreendam que as informações geradas por IA podem conter ``alucinações algorítmicas'' e vieses. A checagem de fatos e a busca por fontes primárias confiáveis são competências essenciais e não negociáveis — especialmente em aplicações de automação industrial, em que decisões de projeto envolvem segurança de pessoas e processos.

    \item \textbf{Transparência e Autoria:} Em conformidade com o princípio da integridade acadêmica, é obrigatório explicitar, de forma clara e transparente, toda e qualquer utilização de ferramentas de IA no processo de concepção, desenvolvimento ou redação de trabalhos acadêmicos. A forma de citação e declaração de uso será orientada pelos docentes em cada componente curricular.

    \item \textbf{Privacidade e Segurança:} É vedada a inserção de dados pessoais sensíveis, informações estratégicas da universidade, de plantas industriais parceiras ou dados de terceiros em plataformas de IA públicas que não garantam os requisitos adequados de segurança e proteção jurídica, em observância à Lei Geral de Proteção de Dados (LGPD).
\end{enumerate}

\subsection{O Papel da IA como Ferramenta de Aprendizagem}

Este curso incentiva o uso consciente da IA como uma poderosa aliada para potencializar a aprendizagem e a criatividade. Encorajamos os(as) estudantes a enxergarem a IA não como um oráculo, mas como um colega de equipe incansável, que pode assumir múltiplos papéis no processo de estudo.

\begin{enumerate}
    \item \textbf{A IA como ``Monitor de Disciplinas em Tempo Integral'':} Os(as) discentes podem utilizar a IA como um assistente de estudos personalizado, disponível a qualquer momento para:
    \begin{itemize}
        \item Explicar conceitos complexos com diferentes analogias.

        \item Sugerir roteiros de estudo e materiais complementares.

        \item Gerar problemas e exercícios para praticar e fixar o conhecimento.

        \item Auxiliar na depuração de programas de CLP, scripts de automação e na compreensão de erros.
        Contudo, reforça-se que este ``monitor'' é um ponto de partida. A validação final do conhecimento e o aprofundamento conceitual devem ser sempre realizados com os docentes e os materiais de referência da disciplina.
    \end{itemize}

    \item \textbf{Combatendo o ``Viés da Primeira Ideia'':} A criatividade, em sua essência, é ``fazer mais do que a primeira coisa que se pensa''~\cite{utley_2025}. A IA é uma ferramenta excepcional para superar a tendência humana de se fixar na primeira solução aparente. Incentiva-se o uso da IA para:
    \begin{itemize}
        \item Gerar um grande volume de ideias e abordagens alternativas para um problema de automação (\textit{brainstorming}).

        \item Explorar diferentes perspectivas e refutar hipóteses iniciais de projeto.

        \item Ampliar o repertório criativo, forçando o(a) estudante a ir além do óbvio.
    \end{itemize}

    \item \textbf{A Importância do Contexto e da Curadoria:} A qualidade da interação com a IA é diretamente proporcional à qualidade do contexto fornecido pelo usuário. Para obter resultados valiosos, o(a) estudante deve atuar como um(a) curador(a) da cognição, fornecendo \textit{briefings} detalhados, questionando as respostas, refinando os comandos (\textit{prompts}) e, principalmente, adicionando sua visão crítica e seu conhecimento de mundo ao produto final. A IA amplifica a capacidade do usuário; ela não a cria do zero~\cite{utley_2025}.
\end{enumerate}

\subsection{Formação e Letramento em IA}

O curso se compromete a promover o letramento em IA de forma contínua e transversal. A disciplina de Introdução à Área Básica de Ingresso em Automação, Robótica e IA (\texttt{FEELT!IABI}) abordará formalmente estas diretrizes, preparando os ingressantes para uma jornada acadêmica consciente. Adicionalmente, os docentes irão, em seus respectivos componentes curriculares — em especial nos conteúdos de inteligência artificial aplicada à automação (\ppcenfasecurricular) —, orientar as boas práticas e os limites do uso de IA, fomentando uma postura de aprendizado e adaptação contínuos, essencial para a atuação profissional nesta nova era.
